\subsection{Prior set-up}\label{app:prior}
For the prior setup, we closely follow \cite{hauzenberger2020effectiveness}. Our priors are designed to provide sufficient regularization of the VAR coefficients to reduce estimation uncertainty. On the VAR coefficients, we therefore apply a Horseshoe (HS) prior \citep{carvalho2010horseshoe}, which belongs to the class of global-local shrinkage techniques. To outline the mechanism of this prior, it proves useful to collect the state-specific coefficients in an $M \times N$-dimensional matrix $\bm A_k = (\bm c_k, \bm \gamma_k, \bm A_kp, \dots, \bm A_kP)$, for $k = \{0, 1\}$ and $p=1,...,P$. In the following, let $\bm a_k = \text{vec}(\bm A_k)$ be a $K (= NM)$-dimensional vector, then, for $k = \{0, 1\}$, then the hierarchical HS prior on the $j$th ($j = 1, \dots, K)$ state-specific coefficients is as follows:
\begin{equation*}
a_{kj} \sim \mathcal{N}(0, \tau_k \lambda_{kj}), \quad \sqrt{\tau_k} \sim \mathcal{C}^+(0,1) \quad \text{and} \quad \sqrt{\lambda_{kj}} \sim \mathcal{C}^+(0,1).   
\end{equation*}
Here, a global component $\tau_k$ pushes all coefficients in $\bm a_k$ toward zero (i.e., the prior mean), while local scaling parameters $\lambda_{kj}$ are regime- and covariate-specific, providing sufficient flexibility.  Both hyperparameters $\sqrt{\tau_{k}}$ and $\sqrt{\lambda_{kj}}$ follow a half-Cauchy distribution denoted by $\mathcal{C}^+$. This prior is able to introduce sufficient sparsity via the global shrinkage component in a highly parameterized model that would otherwise be prone to overfitting. The local scales still allow for the detection of important individual covariates by mitigating the potential (over)shrinkage of the global parameter. 

The prior specification on the variance-covariance $\bm \Sigma$ is standard. Here, $\bm \Sigma$ follows an inverted Wishart distribution: 
\begin{equation*}
    \bm \Sigma \sim \mathcal{IW}(d_0, \bm D_0), 
\end{equation*}
with the scalar $d_0$ denoting the prior degrees of freedom and the $M\times M$-dimensional matrix $\bm D_0$ referring to the prior scaling. In the empirical application, we specify $d_0 = M+2$ and $\bm D_0 = 0.01 \bm I_M$. This choice ensures that the prior is proper, but still weakly informative.

To complete our proposed setup, we also need to specify priors on the hyperparameters $\mu$ and $\rho$, which define the transition between the two regimes. The threshold parameter $\mu$ follows a weakly informative Uniform distribution:
\begin{equation*}
    \rev{\mu \sim \mathcal{U}\left(\bar{z}^* - 0.5,\; \bar{z}^* + 0.5\right),}
\end{equation*}
\rev{with the prior support centered on the sample mean $\bar z^*$ of the (standardised) signal variable, i.e., spanning half a standard deviation on either side of the mean. Widening the support to the full range of $z_t^*$ leaves the posterior of $\mu$ virtually unchanged.}
On the speed of adjustment parameter $\rho$ we impose a Gamma prior: 
\begin{equation*}
    \rho \sim \mathcal{G}(s_0, s_1), 
\end{equation*}
with the hyperparameters \rev{$s_0=4/0.001$ and $s_1=1/0.001$ (prior mean $4$, prior standard deviation $0.06$)} chosen such that the transition between the two states occurs relatively fast. This informative choice ensures a sensible state allocation and a clear separation between the two regimes, one characterized by an expansionary US monetary policy and the other characterized by a restrictive stance (see also \autoref{fig:stallo} for the estimated state allocation). \rev{Centering the prior at $7$ instead of $4$ leaves the state allocation and all impulse responses virtually unchanged; the transition speed is deliberately calibrated, as the likelihood is uninformative about $\rho$.}

\subsection{Posterior quantities and sampling algorithm}\label{app:mcmc}
Next, we outline the main steps of our Markov Chain Monte Carlo (MCMC) algorithm, which is used to simulate from the full posterior distribution. To obtain the conditional posteriors, we combine the priors with the likelihood of the model, as described in \autoref{eq:STVAR} and \autoref{eq:St}. For most parameters, this leads to well-known conditional distributions, and simulating from them is straightforward. Let the symbol $\bullet$ indicate that we condition on the remaining parameters of the system, then these sampling steps read as follows: 
\begin{enumerate}
    \item For simplicity we collect all VAR coefficients in a $2K$-dimensional vector $\bm a = (\bm a_1', \bm a_0')'$ and define the corresponding $2K$-dimensional vector of predictors $\bm x_t = \left((1, z_t, \bm y_{t-1}', \dots, \bm y_{t-p}')\times S_t,  (1, z_t, \bm y_{t-1}', \dots, \bm y_{t-p}') \times (1-S_t)\right)'$. 
    In what follows, we sample $\bm a$ from a multivariate Gaussian distribution: 
    \begin{equation*}
        \bm a|\bullet \sim \mathcal{N}\left(\bar{\bm a}, \bar{\bm V} \right),
    \end{equation*}
    with $\bar{\bm a}$ denoting the posterior mean and $\bar{\bm V}$ the posterior variance-covariance matrix: 
    \begin{equation*}
    \begin{aligned}
       \bar{\bm V} =& \left(\bm \Sigma^{-1} \otimes \bm X'\bm X + \underline{\bm V}^{-1} \right), \\
       \bar{\bm a} =&  \bar{\bm V}^{-1}  \left(\bm \Sigma^{-1} \otimes \bm X'\text{vec}(\bm Y) \right). \end{aligned}
    \end{equation*}
    Here, $\bm X = (\bm x_1, \dots, \bm x_T)'$ and $\bm Y = (\bm y_1, \dots, \bm y_T)'$ denote the full data matrices and $\underline{\bm V} = \text{diag}\left(\{v_{1j}\}_{j=1}^{K}, \{v_{0j}\}_{j=1}^{K}\right)$, where $v_{kj} = \tau_k \lambda_{kj}$, collects the prior variances of all coefficients on the main diagonal. 
    \item To draw the HS scaling parameters we rely on the following steps, outlined in \cite{makalic2015simple}: 
    \begin{equation*}
    \begin{aligned}
        \tau_k|\bullet &\sim \mathcal{IG}\left(\frac{K + 1}{2}, \frac{1}{\phi_k} + \sum_{j = 1}^{K} \frac{a_{kj}}{2\lambda_{kj}}\right), \\
        \lambda_{kj}|\bullet &\sim \mathcal{IG}\left(1, \frac{1}{\zeta_{kj}}+ \frac{a_{kj}}{2 \tau_k}\right), \\
    \end{aligned} 
    \end{equation*}
    where the two auxiliary variables $\phi_k$ and $\zeta_{kj}$ are also drawn from inverted Gamma distributions: 
    \begin{equation*}
    \begin{aligned}
        \phi_k|\bullet &\sim \mathcal{IG}\left(1,1+\frac{1}{\tau_k} \right), \\
        \zeta_{kj}|\bullet &\sim \mathcal{IG}\left(1,1+\frac{1}{\lambda_{kj}} \right).
    \end{aligned} 
    \end{equation*}
    \item We sample $\bm \Sigma$ from an inverted Wishart distribution: 
    \begin{equation*}
        \bm \Sigma|\bullet \sim \mathcal{IW}(d_1, \bm D_1),
    \end{equation*} 
    with posterior degrees of freedom $d_1 = d_0 + T/2$ and $M \times M$-dimensional posterior scaling matrix $\bm D_1 = \bm D_0 + (\bm Y - \bm X \bm a)'(\bm Y - \bm X \bm a)/2 $.
    \item The parameters of the state transition function are sampled jointly with a Metropolis-within-Gibbs step. The candidate values are sampled from independent Gaussian proposal densities: $\mu^{(*)} \sim \mathcal{N}(\mu^{(s)}, \kappa_{\mu})$ and $\rho^{(*)} \sim \mathcal{N}(\rho^{(s)},  \kappa_{\rho})$, with the candidate value $\mu^{(*)}$ ($\rho^{(*)}$) being centered on the last accepted draw $\mu^{(s)}$ ($\rho^{(s)}$) and featuring a variance $\kappa_{\mu}$ ($\kappa_{\rho}$). The acceptance probability of the proposed values is then given by: 
    \begin{equation*}
        \min \left(1, \frac{\mathcal{L}(\mu^{(*)}, \rho^{(*)}|\bullet) p(\mu^{(*)}, \rho^{(*)})}{\mathcal{L}(\mu^{(s)}, \rho^{(s)}|\bullet) p(\mu^{(s)}, \rho^{(s)})} \right),
    \end{equation*}
    with $\mathcal{L}$ denoting the data likelihood conditional on the proposed values ($\mu^{(*)}, \rho^{(*)}$) and last accepted draw ($\mu^{(s)}, \rho^{(s)}$), while $p(\mu^{(*)}, \rho^{(*)})$ and $p(\mu^{(s)}, \rho^{(s)})$ refer to joint prior distribution evaluated at the proposed values and the last accepted draw.
    \rev{In practice, we evaluate $\mathcal{L}$ exactly (up to a proportionality constant that does not depend on $\mu$ and $\rho$) by whitening the system with the error precision matrix. Let $\bm L$ denote the lower Cholesky factor of $\bm\Sigma^{-1} = \bm L \bm L'$ and let $\bm U(\mu, \rho) = \bm Y - \bm X(\mu, \rho)\, \bm A$ collect the reduced-form residuals implied by the candidate transition parameters. The rows of the transformed residual matrix $\bm U(\mu, \rho)\, \bm L$ are standard Gaussian under the model, so that $\log \mathcal{L}(\mu, \rho \mid \bullet) = -\frac{1}{2}\lVert \bm U(\mu, \rho)\, \bm L \rVert_F^2 + \text{const}$, where $\lVert \cdot \rVert_F$ is the Frobenius norm. This makes the Metropolis acceptance ratio an exact evaluation of the conditional posterior of $(\mu, \rho)$.}\footnote{\rev{Note for revision (not for publication): an earlier implementation whitened with the transposed factor $\bm L'$ instead of $\bm L$, which leaves the transformed residual covariance different from the identity matrix and hence renders the acceptance ratio inexact. Correcting this (code: \texttt{STVAR\_rev\_whitenfix.R}, July 2026) leaves all impulse responses and the state allocation virtually unchanged (posterior-median IRF correlations $\geq 0.997$; regime-path correlation $0.99$), but changes the posterior of $\rho$: under the exact likelihood the posterior of $\rho$ essentially coincides with its informative prior, i.e., the speed of adjustment is pinned down by the prior rather than the likelihood, while $\mu$ remains data-identified. Metropolis acceptance rises from $0.09$ to $0.63$.}}
\end{enumerate}
\rev{We repeat these steps $12,000$ times and discard the initial $2,000$ draws as burn-in, retaining $10,000$ draws for inference. [Note for revision: iteration count and the prior settings for $\rho$ and $\mu$ above have been harmonized with the replication code (previously the text stated $15,000/5,000$ iterations, a $\rho$-prior mean of $7$, and min--max support for $\mu$). A check re-estimating the baseline under the text's former settings leaves all results virtually unchanged: regime-path correlation $0.995$, IRF-median correlations $\geq 0.997$, identical hawk--dove contrast.]}